\scanpage{348}
% top centre
\mspage{264}
% lines 1–2
$(\unsure{\mathcal E|\mathscr E|\mathcal C}_{\unsure{j|\pi}})_\pi$ désignant la sous-catégorie pleine de $\mathcal E_j$ formée des objets $\pi$-accessibles de $\mathcal E_j$.

% lines 3–12
En effet, \unsure{les'|les|l'} hypothèses \unsure{du début|dites|essentielles} impliquent l'existence de $\pi_0$ et $\pi$ satisfaisant les conditions dites. Les conditions a) et b) sur $\pi$, \deleted{$\unsure{\pi}$} plus le fait que les foncteurs $\lambda_*$ soient $\pi_0$-accessibles, impliquent que la sous-catégorie \unsure{str.} pleine $N_\pi$ de $N$ \deleted{formée des objets}\inserted{dont les fibres} sont les $(N_j)_\pi$, est une sous-catégorie cofibrée, et que de plus $N$ est $J$-équivalente à $\operatorname{Ind}_\pi(N_\pi/J)$.

% lines 13–18
\deleted{La condition} D'autre part, les fibres de $N_\pi/J$ sont stables par $\varinjlim$ \unsure{grdes|gdes} devant $\pi_0$ et de cardinal $\leq\pi$, donc la condition ``\unsure{ordinale|ordinaire}'' de 4.10.1. est vérifiée \inserted{(en y faisant $D=N_\pi$)}. \deleted{\unsure{D}} Comme de plus $\mathrm{Card}\,J\leq\pi$, on en conclut \inserted{par 4.10.1} que \deleted{$\unsure{N\simeq|M\simeq}$} $M\approx\operatorname{Ind}_\pi(M_\pi)$.
% archive sheet number: 347
